
The Lower Hybrid current drive problem is one the more important addressed
by Fokker-Planck calculations in the field of fusion by magnetic
confinement. It is also certainly one of the most demanding case regarding
the numerical constraints that must be satisfied simultaneously such as
possible discontinuities and large values of the quasilinear diffusion
coefficient, presence of cross-derivatives $D_{p\xi }=D_{\xi p}\neq 0.$ In
that conditions, the matrix conditionning may be significantly reduced,
leading to possible onset of numerical instabilities, lack of convergence or
possible convergence towards a solution that has no physical sense. In these
conditions, the demonstration that the code remains fully conservative and
gives a physical solution for a domain of parameters that is relevant for
magnetic fusion simulations is a stringent test for the validity of the
projection technique of the kinetic equations on the numerical grids as well
as the correct description of both the internal and external boundary
conditions.

For the test procedure, a simplified expression of the quasilinear diffusion
operator is considered, as described in Appendix \ref{simplifiedlhcd}, using
the familiar boxcar shape in momentum space $\overline{D}_{0,new}^{LH}=1$
for $v_{1}\leq v_{\Vert }\leq v_{2},$ and $\overline{D}_{0,new}^{LH}=0$
otherwise, neglecting the factor $v_{th}^{\dag }/v_{\Vert },$ as done
usually in the litterature which only acceptable in principle when $%
\overline{D}_{0,new}^{LH}\gg 1$ (see Refs. \cite{kar85},\cite{kar86},\cite
{sho94},\cite{pey98}). For comparison with existing results, simulations are
performed for the couples of values $\left( v_{1}=3.5,v_{2}=6\right) $ and $%
\left( v_{1}=4,v_{2}=7\right) .$ Except it is specified, the default upper
bound of the integration domain is taken at $p_{\max }=20$. Plasma
parameters are similar to those used in Sec.\ref{conducsection} dedicated to
the the Ohmic conductivity problem, except that the reference electron
temperature may take two different values $T_{e}^{\dagger }=5.11\times
10^{-4}keV$ and $T_{e}^{\dagger }=5.11keV$ corresponding to non-relativistic
to relativistic limits respectively. The choice of these electron
temperatures allows comparaison with already published results at $\beta
_{th}^{\dag }=0.01\left( \approx 0\right) $ and $\beta _{th}^{\dag }=0.1$.
All plasma parameters are gathered in the file ``CQL3D\_LH'', since it
corresponds to conditions that have been used for code comparison with the
well known CQL3D code \cite{har93c1_IAEA}.

Concerning the numerical parameters, the code is running by default in the
fully implicit mode $\Delta \overline{t}=10000,$ with the linearized
electron-electron Belaiev-Budker collision operator that conserve momentum.
The drop tolerance for incomplete $LU$ matrix factorization is $10^{-4}$,
while no Maxwellian solution is imposed at $i=1/2$. The momentum and
pitch-angle grids are taken uniform with $n_{p}=n_{\xi _{0}}\simeq 200$
except when simulations with different numbers of grid points or non-uniform
meshes are studied. Furthermore, no bounce averaging is considered, expect
if specified.

The first problem that is addressed is the numerical accuracy with the
number of grid points. For a pure hydrogen plasma $\left( Z_{s}=1\right) ,$ $%
n_{p}=n_{\xi _{0}}$ are varying from $50$ to $250$. Here a very simple case
is considered, with $\left( v_{1}=3,v_{2}=5\right) ,$ using the relativistic
Maxwellian collision operator, but at $\beta _{th}^{\dag }=0.001$
(non-relativistic limit). This allows easy comparison with previously
published results \cite{kar86}. As shown in Fig.\ref{efflhgrille},
significant oscillations of the absorbed RF $\left\langle
P_{LH}\right\rangle $ power density arises when $\left( n_{p},n_{\xi
_{0}}\right) \leq 150$ which correspond to the relative position of the
lower bound of the Lower Hybrid resonance domain with respect to the nearest
grid point. The exponential decrease makes the solution very sensitive to
this parameter. For this domain of parameter, the power density absorbed by
collisions $\left\langle P_{coll}\right\rangle $ is not exactly equivalent
to $\left\langle P_{LH}\right\rangle ,$ though close to unity within $5\%$,
leading to oscillations of the ratio $\left\langle P_{coll}\right\rangle
/\left\langle P_{LH}\right\rangle .$ The current density is clearly less
sensitive to the grid size, since it is mainly driven by the upper part of
the tail. Nevertheless, small oscillations are observed for a coarse grid.
Consequently, the current drive efficiency is varying significantly when $%
\left( n_{p},n_{\xi _{0}}\right) \leq 150,$ as shown in Fig. \ref
{efflhgrille}, and a reltive accuracy lower than $0.5\%$ is only found when $%
\left( n_{p},n_{\xi _{0}}\right) \geq 150.$ From this systematic study, is
turns out that accurate calculations with uniform pitch-angle and momentum
grids require discrete steps lower than $\left( \Delta p,\Delta \xi
_{0}\right) \lesssim 0.13$. The asymptotic value of the current drive
efficiency is found to reach $14.26$, very close to the value $14.35$ given
in Ref. \cite{kar86}, despite the fact that $\left\langle
P_{LH}\right\rangle $ and $\left\langle j_{LH}\right\rangle $ found by the
code are larger by more than $10\%$. The difference arises likely from the
collision operator that may exhibit some differences.

%TCIMACRO{
%\TeXButton{Fig_efflhgrille}{\fig{figures/efflhgrille.eps}{Variations of the Lower Hybrid current and power densities, ratio between the RF and collision absorbed power density, 
%and the current drive efficiency with the grid size. Here uniform pitch-angle and momentum grids are considered. Detailed aspect of the simulation 
%are given in the text}{efflhgrille}{7cm} }}
%BeginExpansion
\fig{figures/efflhgrille.eps}{Variations of the Lower Hybrid current and power densities, ratio between the RF and collision absorbed power density, 
and the current drive efficiency with the grid size. Here uniform pitch-angle and momentum grids are considered. Detailed aspect of the simulation 
are given in the text}{efflhgrille}{7cm} %
%EndExpansion

The effect of the grid size of memory storage requirements and the full
elapsed time for kinetic calculations. A shown in Fig.\ref{memlhgrille} the
memory used for data storage of $\widehat{\Bbb{L}}$ and $\widehat{\Bbb{U}}$
grows quadratically as expected from a $2-D$ problem. The dashed line is a
parabolic fit that confirms this dependence. However, even if the scaling is
quadratic, memory requirement remains reasonable, because of the small
coefficients pruning in the matric factorization procedure. In the case here
considered, the drop tolerance is fairly low, $10^{-4},$ but large values
may drastically reduce this size as discussed in Sec.\ref{LUsection}. Much
in the same way, the full elapsed time for kinetic calculations increases
also quadratically, but since time step is very large, is remains
nevertheless lower than one minute. This confirms the effectiveness of the
code, which is fast, accurate and robust, therefore fully designed for
self-consistent realistic calculations that involve a chain of codes.

%TCIMACRO{
%\TeXButton{Fig_memlhgrille}{\fig{figures/memlhgrille.eps}{Variations of the memory storage requirement and the time elapsed for kinetic calculations with the grid size. 
%Here uniform pitch-angle and momentum grids are considered. Detailed aspect of the simulation are given in the text}{memlhgrille}{7cm} }}
%BeginExpansion
\fig{figures/memlhgrille.eps}{Variations of the memory storage requirement and the time elapsed for kinetic calculations with the grid size. 
Here uniform pitch-angle and momentum grids are considered. Detailed aspect of the simulation are given in the text}{memlhgrille}{7cm} %
%EndExpansion

For a fixed number of grid points, the upper limit of the integration domain
may have a crucial role, especially when entering relativistic regimes at $%
\beta _{th}^{\dag }=0.1$ and when the resonance domain involves large values
of $v_{\Vert }$. As shown in Fig. \ref{efflhpmax} for a pure hydrogen plasma 
$\left( Z_{s}=1\right) ,$ $p_{\max }$ must be larger than $25$ for $\left(
v_{1}=4,v_{2}=7\right) $ so that $\eta _{LH}$ becomes independent of $%
p_{\max },$ while $p_{\max }\simeq 15$ is acceptable for $\left(
v_{1}=3.5,v_{2}=6\right) .$ The ratio between the power absorbed by the
Lower Hybrid wave and by collisions is also a good estimate of the
robustness of the calculations regarding this parameter. In principle, this
ratio must be close to unity, but significant departures are observed when
the upper limit of the integration domain in momentum space is too low.
Consequently, for robust estimate of the current drive efficiency in
realistic tokamak simulations, $p_{\max }=30$ must be considered as a
reference value for relativistic calculations, providing the lower bound of
the resonance domain is larger than $n_{\Vert }\simeq 1.4.$ As shown in Sec. 
\ref{boundarylhsec}, values of $n_{\Vert }$ lower than $1.4$ becomes rapidly
highly questionable when $T_{e}^{\dagger }$ exceeds $1keV$, because of the
relativistic curvature of the resonance domain that makes the code
fondamentaly non conservative, like in a runaway regime. Considering that in
most plasma conditions encountered in tokamaks, $n_{\Vert }\geq 1.4,$ the
choice of $p_{\max }=30$ makes the code fairly insensitive to the
characteristics of the Lower Hybrid wave spectrum. This is an important
aspect for the robustness of self-consistent simulations. Since the upper
value of the momentum domain of integration is fairly large, the main
consequence is that the number of mometum or pitch-angle grid points must
exceed $230$.

%TCIMACRO{
%\TeXButton{Fig_efflhpmax}{\fig{figures/efflhpmax.eps}{Variation of the current drive efficiency with the upper mometum limit of the integration domain}{efflhpmax}{7cm} }}
%BeginExpansion
\fig{figures/efflhpmax.eps}{Variation of the current drive efficiency with the upper mometum limit of the integration domain}{efflhpmax}{7cm} %
%EndExpansion

The physical benchmark of the code is performed by investigating the role
played by increasingly large values of $Z_{s}$. Though this corresponds to
unrealistic situations, it allows interesting comparisons with theoretical
results. As shown in Fig. \ref{efflhZi}, the current drive efficiency
defined as the ratio 
\begin{equation}
\eta _{LH}=\frac{\left\langle j_{LH}\right\rangle }{\left\langle
P_{LH}\right\rangle }
\end{equation}
is decreasing as $Z_{s}$ rises as expected because of enhanced pitch-angle
scattering, either for $\left( v_{1}=3.5,v_{2}=6\right) $ or $\left(
v_{1}=4,v_{2}=7\right) .$ Relativistic effects are almost negligible for the
case $\left( v_{1}=3.5,v_{2}=6\right) $ but also for $\left(
v_{1}=4,v_{2}=7\right) $ when the upper limit of the momentum integration
domain is $p_{\max }=30.$ For $p_{\max }=20,$ the current drive efficiency
is significantly underestimated, as discussed above, when large values of $%
v_{\Vert }$ are considered, and when relativistic effects become
significant, i.e. when $\beta _{th}^{\dag }=0.1$ in the case here
considered. The current drive efficiency scales roughly with the $1-D$
non-relativistic model given in Ref. \cite{kar86} 
\begin{equation}
\eta _{LH}^{1-D}\simeq 0.7\frac{4}{5+Z_{s}}\frac{v_{2}^{2}-v_{1}^{2}}{\ln
\left( v_{2}/v_{1}\right) }
\end{equation}
where the dependence $4/\left( 5+Z_{s}\right) $ with $Z_{s}$ is deduced from
Langevin analysis. It is found that the coefficient $0.7$ is independent of $%
\left( v_{1},v_{2}\right) $. This factor could arise from that fact that the 
$1-D$ expression is obtained in the saturated limit corresponding to $%
\overline{D}_{0,new}^{LH}=+\infty ,$ while numerical calculations are
obtained with $\overline{D}_{0,new}^{LH}=1.$ This result confirms on one
side the role played by pitch-angle scattering that make the kinetic problem
a $2-D$ one, and that on the other side, important aspects of the $1-D$
description remains still valid. The $1-D$ aspect of the Lower Hybrid
physics is supported by the ratio $\eta _{LH}\left[ v_{1}=3.5,v_{2}=6\right]
/\eta _{LH}\left[ v_{1}=4,v_{2}=7\right] $ which is close to the expected
theoretical value $0.75$ when $Z_{s}$ is not too large. In this regime,
electron-electron interactions deeply contribute to the current drive
process that is fundamentaly 1-D in momentum space, since resonace
acceleration takes place along the magnetic field line. When $Z_{s}\gg 1,$
the ratio of current drive efficiencies drop much faster, since $2-D$
effects start to become important. From these tests, it is clear that the
code captures most of the features of the Lower Hybrid current drive. The
parametric dependence with $Z_{s}$ indicates that $1-D$ physics predominates
over the $2-D$ one provided $Z_{s}$ is not too large. The $2-D$ effects are
just corrections that may be incorporated in a single reduction factor. It
is important to mention that this result is independent of the model for
describing electron-ion collisions, considering an ion Maxwellian
background, or the high-velocity limit.

%TCIMACRO{
%\TeXButton{Fig_efflhZi}{\fig{figures/efflhZi.eps}{Variation of the current drive efficiency with the main ion charge in the plasma}{efflhZi}{7cm} }}
%BeginExpansion
\fig{figures/efflhZi.eps}{Variation of the current drive efficiency with the main ion charge in the plasma}{efflhZi}{7cm} %
%EndExpansion

Successful comparisons between different Fokker-Planck codes have been
performed. Results are given in Table \ref{table_lh1} for a uniform grid,

\begin{table}[tbp]
\begin{center}
\begin{tabular}{|l|l|l|}
\hline
$\left( v_{1},v_{2}\right) $ & $\eta _{LH}$ [DKE code] & $\eta _{LH}$ \cite
{sho94} & $\eta _{LH}$ \cite{kar86} \\ \hline
$\left( 3.5,6\right) $ & $21.4$ & $22.3$ & $-$ \\ \hline
$\left( 4,7\right) $ & $28.38$ & $28.8$ & $29.7$ \\ \hline
\end{tabular}
\end{center}
\caption{Lower Hybrid current drive efficiencies in a pure hydrogen plasma
from various $2-D$ relativistic Fokker-Planck codes}
\label{table_lh1}
\end{table}
considering the current efficiency in normalized units, but also an
excellent agreement is found in absolute units, as shown in Table \ref
{table_lh2}
\begin{table}[tbp]
\begin{center}
\begin{tabular}{|l|l|l|}
\hline
$\left( v_{1},v_{2}\right) $ & $\eta _{LH}$ [DKE code] & $\eta _{LH}$ \cite
{kar79} & $\eta _{LH}$ \cite{har93c1_IAEA} \\ \hline
$\left( 4,7\right) $ & $32$ & $31.96$ & $31.34$ \\ \hline
\end{tabular}
\end{center}
\caption{Lower Hybrid current drive efficiencies $\left( A\cdot m/W\right) $
in a pure hydrogen plasma from various $2-D$ relativistic Fokker-Planck codes
}
\label{table_lh1}
\end{table}
taking the same value for the Coulomb logarithm, $\ln \Lambda \approx 15.$

For a uniform grid in $p$ and $\xi _{0}$, $\overline{D}_{0,new}^{LH}$ may be
increased significantly above unity, before onset of a spurious numerical
behaviour. Simulation $\beta _{th}^{\dag }=0.1$ have been runned successfuly
with $\overline{D}_{0,new}^{LH}=10,$ keeping a fully conservative code,
while the criterion $\left\langle P_{coll}\right\rangle /\left\langle
P_{LH}\right\rangle \approx 1$ is satisfied, which validates the physical
nature of the solution found numericaly. This gives quite confortable
margins for coupling the kinetic code within ray-tracing calculations in the
weak absorption regime where $\overline{D}_{0,new}^{LH}$ my be large.
However, there are natural limitations for $\overline{D}_{0,new}^{LH},$ that
arise from the discrete projection of the differential equation on the
numerical grid. Above $\overline{D}_{0,new}^{LH}=1,$ there is a risk that
that code becomes unstable. In that case, typical signs are the lack of
particle conservation, and the wrong balance between $\left\langle
P_{coll}\right\rangle $ and $\left\langle P_{LH}\right\rangle .$ Several
strategies may be considered to overcome this difficulty. As a first
attempt, the reduction of $\Delta \overline{t}$ improves the matrix
conditionning, since the main diagonal becomes more prominant over
off-diagonal coefficients. If this approach fails, increasing the grid size
may help to fix this problem in some limits. In any case, forcing the
Maxwellian solution close to $p=0$ and normalizing the density at each
iteration in this kind of regime is usualy useless, since the numerical
solution may reamin fully non physical, though convergence is ensured. In
that case, the overall absorption process fails, and the result has no
meaning. A possible way to overcome this problem is to perform calculations
with the upper acceptable $\overline{D}_{0,new}^{LH}$ value using an adhoc
correcting factor, in order to recover the solution corresponding asymptotic
limit at $\overline{D}_{0,new}^{LH}=+\infty .$ 

When a non-uniform pitch-angle grid is used, most of the results that are
obtained are similar, provided $\overline{D}_{0,new}^{LH}$ is close to $1$,
and $Z_{s}\leq 2.$ For larger values of both parameters, the current drive
efficiency is found to drop by a signficant factor. 
